\documentclass[twocolumn]{aastex631}
\begin{document}
\title{A residual reionization ionizing-photon-budget shortfall for star-forming galaxies at z\~{}9 (jwsthigh systematic corner)}
\author{NebulaMind Lab (autonomous pipeline)}
\affiliation{NebulaMind Lab, \url{https://lab.nebulamind.net}}
\begin{abstract}
Reionization ionizing-photon-budget reconciliation at z\~{}9: star-forming galaxies require f\_esc=0.208 (+0.197/-0.100) to close the budget (Madau-Dickinson SFRD, log xi\_ion=25.5+/-0.15, clumping C=2-5, JWST-SFRD tail), versus indirect-proxy-inferred f\_esc=0.062 (+0.110/-0.039) from LzLCS O32/beta calibrations. Median delta(required-inferred)=+0.129 dex-frac (16-84\%: +0.000 to +0.328); 84\% of the systematic MC shows a shortfall. Conclusion: a genuine SHORTFALL remains, and the sign holds under both O32 and beta calibrations. The result is bounded by the xi\_ion x clumping x proxy-calibration systematic, not by statistics. Generated autonomously from public data (jwst) via the NebulaMind Lab runner: a bounded, reproducible, descriptive study.
\end{abstract}
\section{Introduction}
Recent studies have highlighted a potential crisis in our understanding of reionization, with concerns that current models may not account for the necessary ionizing photons to drive this process [Muñoz2024]. This issue is further complicated by the demands placed on ionizing sources during absorption-dominated reionization [Davies2021] and the need for accurate calibration of excursion set reionization models [Park2022]. To address these challenges, it is essential to reassess the galaxy ionizing photon budget at high redshifts [Duncan2015].
\section{Data and method}
In this work, we adopt a literature-anchored approach to calculate the reionization-photon-budget. We utilize the Madau \& Dickinson (2014) analytic fitting function for the cosmic star formation rate density (SFRD), along with published values for xi\_ion and the O32/beta f\_esc proxy calibrations from LzLCS [Chisholm+22, Flury+22; Simmonds+24]. Our method focuses on reconciling systematic uncertainties in ionizing-photon-budget calculations using a Monte Carlo approach to quantify the required escape fraction (f\_esc) of ionizing photons.
\section{Result}
Our analysis reveals that at z\~{}9, star-forming galaxies must have an f\_esc value of 0.208 (+0.197/-0.100) to reconcile the reionization photon budget under the Madau-Dickinson SFRD, log xi\_ion=25.5+/-0.15, and clumping factor C=2-5. However, indirect-proxy-inferred f\_esc values from LzLCS O32/beta calibrations yield a significantly lower estimate of 0.062 (+0.110/-0.039). The median difference between the required and inferred escape fractions is +0.129 dex-frac (16-84\%: +0.000 to +0.328), with 84\% of our systematic Monte Carlo simulations indicating a shortfall in ionizing photons.
\begin{figure}\centering\includegraphics[width=\columnwidth]{result.png}\caption{A residual reionization ionizing-photon-budget shortfall for star-forming galaxies at z\~{}9 (jwsthigh systematic corner)}\end{figure}
\section{Caveats}
It is crucial to acknowledge the limitations of our approach, which relies on an automated, single-selection, uncalibrated measurement. This method may not fully capture the complexities of reionization and the escape fraction of ionizing photons. Additionally, our analysis depends on the accuracy of previously published values for xi\_ion and O32/beta f\_esc proxy calibrations, which could introduce uncertainties into our results. Furthermore, the use of a single SFRD model (Madau \& Dickinson 2014) may not account for variations in star formation rates across different galaxy populations. These caveats highlight the need for further research and refined models to better understand the reionization process and its underlying mechanisms.
\section*{References}
\footnotesize
[Muoz2024] Muñoz et al. (2024). Reionization after JWST: a photon budget crisis?. bibcode 2024MNRAS.535L..37M \par [Davies2021] Davies et al. (2021). The Predicament of Absorption-dominated Reionization: Increased Demands on Ionizing Sources. bibcode 2021ApJ...918L..35D \par [Park2022] Park et al. (2022). Calibrating excursion set reionization models to approximately conserve ionizing photons. bibcode 2022MNRAS.517..192P \par [Duncan2015] Duncan et al. (2015). Powering reionization: assessing the galaxy ionizing photon budget at z < 10. bibcode 2015MNRAS.451.2030D \par [Madau2017] Madau (2017). Cosmic Reionization after Planck and before JWST: An Analytic Approach. bibcode 2017ApJ...851...50M

\end{document}
